\begin{table*}[t]
	\centering
	\scriptsize
	
	\begin{tabular}{| l | c || c | c | c | c | c || c | c | }
	\hline
	{\bf{Environment}} &{\bf{Ours}} &
	{\bf{Bisimulation}} &
	{\bf{Branching}} &
	{\bf{Boltzmann}} &
	{\bf{GNN Sh}} &
	{\bf{GNN MF}} &
	{\bf{Manual}} &
	{\bf{No Invent}} \\
	\hline
	PickPlace1D & 625 (134) & 176 (3) & 219 (145) & 264 (17) & 1951 (85) & 1951 (85) & 177 (147) & 66 (0) \\
	Blocks & 10237 (853) & 800 (44) & 1561 (98) & 9798 (1688) & 4047 (209) &4047 (209) & 102 (5) & 84 (2) \\
	Painting & 18395 (28153) & 872 (380) & 2883 (144) & 9457 (3421) & 9185 (166) & 9185 (166) & 565 (457) & 260 (2) \\
	Tools & 18666 (2815) & 573 (20) & 5524 (747) & 9716 (1000) & 7362 (197) & 7362 (197) & 167 (3) & 141 (3) \\\hline
	\end{tabular}
	\caption{\textbf{Learning times in seconds for all experiments}. All numbers are means over 10 seeds, with standard deviations in parentheses. For the GNN-based methods, learning time encompasses training the neural networks. For the other methods, learning time encompasses learning predicates, operators, and samplers (i.e., all components of the abstraction). Even though our main method performs well (Ours), this does come at the cost of increased learning time (although the learning is purely offline). Note that the Manual approach only manually specifies a \emph{state} abstraction (predicates); operators and samplers must still be learned, contributing to the non-zero learning time. Thus, comparing Ours and Manual shows that the large majority of learning time in our system is spent on predicate invention.}
	\label{tab:learningtimeresults}
\end{table*}